\documentclass[10pt,a4paper]{article}

\usepackage[margin=20mm]{geometry}
\usepackage{amsmath,amssymb,amsthm}

\theoremstyle{definition}
\newtheorem*{definition*}{Definition}
\newtheorem*{problem*}{Problem}

\newcommand{\CSPACE}{\mathsf{CSPACE}}
\newcommand{\CTISP}{\mathsf{CTISP}}
\newcommand{\CL}{\mathsf{CL}}
\newcommand{\CLP}{\mathsf{CLP}}
\newcommand{\Lclass}{\mathsf{L}}
\newcommand{\NL}{\mathsf{NL}}
\newcommand{\TC}{\mathsf{TC}}
\newcommand{\Pclass}{\mathsf{P}}
\newcommand{\ZPP}{\mathsf{ZPP}}
\newcommand{\LOSSY}{\mathsf{LOSSY}}

\setlength{\parindent}{0pt}
\setlength{\parskip}{5pt}
\setlength{\emergencystretch}{2em}

\begin{document}

\begin{center}
  {\Large\bfseries Catalytic Logspace versus Polynomial Time}
\end{center}

\section{Definitions}

\begin{definition*}[Catalytic space and $\CL$]
A \emph{catalytic Turing machine} is a deterministic Turing machine with a
read-only input tape, an ordinary read-write work tape, and an additional
read-write catalytic tape. On input $x$, the catalytic tape initially contains
an arbitrary string $a$. The machine must halt for every pair $(x,a)$, return
the same decision for every choice of $a$, and restore the entire catalytic
tape exactly to $a$ before halting.

For functions $S,S_a:\mathbb N\to\mathbb N$, let
$\CSPACE(S(n),S_a(n))$ denote the class of languages decided using $O(S(n))$
work-tape cells and $O(S_a(n))$ catalytic-tape cells. Write
\begin{equation}
  \CSPACE(S(n))
  :=
  \CSPACE\bigl(S(n),2^{O(S(n))}\bigr)
\end{equation}
and define
\begin{equation}
  \CL:=\CSPACE(\log n).
\end{equation}
Thus a $\CL$ machine has $O(\log n)$ ordinary work space and a polynomial-size
catalyst. Its definition imposes no polynomial running-time bound.
\end{definition*}

\begin{definition*}[Polynomial-time catalytic logspace]
The class
\begin{equation}
  \Pclass
  :=
  \bigcup_{k\geq1}\mathsf{DTIME}(n^k)
\end{equation}
contains the languages decidable by ordinary deterministic Turing machines in
polynomial time. Let $\CTISP(t,s,c)$ denote catalytic computations with time,
ordinary-space, and catalytic-space bounds $O(t)$, $O(s)$, and $O(c)$,
respectively. Define
\begin{equation}
  \CLP
  :=
  \bigcup_{k\geq1}\CTISP(n^k,k\log n,n^k).
\end{equation}
Consequently,
\begin{equation}
  \CLP\subseteq\CL\cap\Pclass.
\end{equation}
\end{definition*}

\begin{definition*}[Lossy coding]
An instance of the total search problem $\mathsf{LossyCode}$ is a pair of
Boolean circuits
\begin{equation}
  C:\{0,1\}^m\longrightarrow\{0,1\}^{m-1},
  \qquad
  D:\{0,1\}^{m-1}\longrightarrow\{0,1\}^m.
\end{equation}
The task is to output a string $z\in\{0,1\}^m$ such that
\begin{equation}
  D(C(z))\neq z.
\end{equation}
Such a string always exists: the map $D\circ C$ factors through a set of only
$2^{m-1}$ codewords and therefore cannot fix all $2^m$ inputs. The class
$\LOSSY$ consists of the languages decidable in deterministic polynomial time
with oracle access to a function that returns a valid solution to every
$\mathsf{LossyCode}$ instance.
\end{definition*}

\newpage
\section{Best Known Result}

The known lower and upper bounds are
\begin{equation}
  \Lclass
  \subseteq
  \NL
  \subseteq
  \TC^1
  \subseteq
  \CLP
  \subseteq
  \CL
  \subseteq
  \LOSSY
  \subseteq
  \ZPP.
\end{equation}
At the level of general complexity classes, $\TC^1$ is the strongest circuit
class stated in the source to be contained in $\CL$. Important individual
problems beyond this circuit-class frontier are known to belong to $\CLP$,
including linear matroid intersection, but this does not yield a larger
standard circuit-class containment.

The best general upper bound is
\begin{equation}
  \CL\subseteq\LOSSY\subseteq\ZPP.
\end{equation}
The second containment follows because a uniformly random
$z\in\{0,1\}^m$ is a valid $\mathsf{LossyCode}$ solution with probability at
least $1/2$, and validity can be checked deterministically. No deterministic
polynomial-time algorithm is known for the general $\mathsf{LossyCode}$
problem.

The structural identity
\begin{equation}
  \CLP=\CL\cap\Pclass
\end{equation}
is known. It follows immediately that
\begin{equation}
  \CL\subseteq\Pclass
  \quad\Longleftrightarrow\quad
  \CL=\CLP.
\end{equation}
Thus proving that every catalytic-logspace language has an ordinary
polynomial-time algorithm automatically gives it a polynomial-time catalytic
algorithm as well.

\section{Problem Formalization}


Prove or disprove the unconditional deterministic containment
\begin{equation}
  \CL\subseteq\Pclass.
\end{equation}

\end{document}
